\section{Función de transferencia}

Se define $f(t)$ como la entrada y $y(t)$ como la salida. Para obtener la función de transferencia se aplica la transformada de Laplace a \eqref{eq:movimiento} bajo condiciones iniciales nulas,
\begin{equation}
  y(0)=0, \qquad \dot y(0)=0.
\end{equation}
En consecuencia,
\begin{equation}
  ms^2Y(s)+\beta sY(s)+kY(s)=F(s).
\end{equation}
Factorizando $Y(s)$ resulta
\begin{equation}
  Y(s)\left(ms^2+\beta s+k\right)=F(s),
\end{equation}
y la función de transferencia entre la fuerza perturbadora y el desplazamiento es
\begin{equation}
  \boxed{G(s)=\frac{Y(s)}{F(s)}
  =\frac{1}{ms^2+\beta s+k}}.
  \label{eq:transferencia}
\end{equation}

Para reconocer la estructura estándar de segundo orden se divide la ecuación de movimiento entre $m$:
\begin{equation}
  \ddot y+\frac{\beta}{m}\dot y+\frac{k}{m}y
  =\frac{1}{m}f(t).
  \label{eq:normalizada}
\end{equation}
Al comparar su lado izquierdo con
\begin{equation}
  \ddot y+2\zeta\omega_n\dot y+\omega_n^2y,
\end{equation}
se identifican
\begin{equation}
  \boxed{\omega_n=\sqrt{\frac{k}{m}}},
  \qquad
  \boxed{\zeta=\frac{\beta}{2\sqrt{mk}}}.
  \label{eq:parametros_dinamicos}
\end{equation}
Por tanto, \eqref{eq:transferencia} también puede escribirse como
\begin{equation}
  G(s)=
  \frac{\dfrac{1}{k}\omega_n^2}
  {s^2+2\zeta\omega_ns+\omega_n^2}.
  \label{eq:transferencia_estandar}
\end{equation}
Esta forma muestra que la ganancia estática es $G(0)=1/k$ y que el denominador determina la frecuencia y el amortiguamiento de la respuesta transitoria.

